\section{Tools for moving base components}
\label{sec:common-tools}

This section follows the base components of a complete kernel of jets
as the marked point moves. Their multiplicities give degree and
canonical-intersection bounds. A separate finiteness lemma treats the
covers that arise when a component projects generically finitely onto
a fixed base. These arguments use neither nonvanishing, abundance,
minimal models, nor logarithmic subadditivity.
Section~\ref{sec:jets} applies them to moving components produced by
failure of the scalar or two-slot jet estimates.

All varieties in the moving-center arguments are over $\C$.
A rational Cartier divisor is a rational divisor some positive
integral multiple of which is Cartier. A requirement that $kP$
be Cartier means that $kP$ is an actual integral Cartier divisor.
For a nef
class on an integral subvariety, intersections are computed after
pullback to its normalization and a resolution. For a linear
subseries of a line bundle, its base ideal is the image of its
evaluation map after tensoring by the inverse line bundle.
``Isolated at the generic point of $V$'' means that this ideal
is primary for the maximal ideal of $\mathcal O_{Z,\eta_V}$.
All powers of ideals below are ordinary powers.

\subsection{Numerical subadjunction on the center}

The first result converts a generic log canonical center into a
canonical-intersection inequality. Its testing class is allowed
to live on the center itself. The proof uses dlt adjunction and
the klt-trivial canonical bundle formula; it does not require the
auxiliary pair to be log canonical away from the generic center.

\begin{lemma}[Numerical generic subadjunction]
\label{common:numerical-subadjunction}
Let $Z$ be a projective $\Q$-factorial klt variety over $\C$, let
$\Theta\geq0$ be a rational divisor, and let $V\subset Z$ be an
integral positive-dimensional subvariety. Suppose that $(Z,\Theta)$
is lc at the generic point of $V$ and that a divisor of log
discrepancy zero has center $V$. Let $f=\dim V$, let $P$ be a nef
rational Cartier divisor on $V$, and let $\rho:V^*\to V$ be a smooth
projective resolution, factoring through the normalization. Then
\begin{equation}
\label{common:subadjunction-inequality}
 K_{V^*}\cdot(\rho^*P)^{f-1}
 \leq (K_Z+\Theta)|_V\cdot P^{f-1}.
\end{equation}
The right side is the intersection on $V$ of the restricted
rational Cartier class with $P^{f-1}$. No $\Q$-Gorenstein hypothesis
on the normalization of $V$ is required.
\end{lemma}

\begin{proof}
Apply the dlt blowup theorem for an effective rational divisor
\cite[Theorem~2.10]{FujinoHashizume2023}. It gives a projective
birational morphism $p:Q\to Z$ with $Q$ $\Q$-factorial and
\begin{equation}
\label{common:subadjunction-surplus}
 p^*(K_Z+\Theta)=K_Q+B_Q+E_Q,
\end{equation}
where $(Q,B_Q)$ is dlt and $E_Q\geq0$ is the surplus over the
non-lc locus. In this construction the coefficients of the strict
boundary are truncated at one and the relevant extracted divisors
are reduced; the coefficients left over form $E_Q$. The surplus
is absent over the open where $(Z,\Theta)$ is lc. In particular it
is absent over a neighborhood of the generic point of $V$.

There are lc strata of $(Q,B_Q)$ mapping onto $V$. Indeed over that
lc neighborhood the construction is crepant, and the center of a
log-discrepancy-zero divisor on the dlt model is an lc stratum.
Choose a stratum $S$ minimal among those mapping onto $V$.
Iterated dlt adjunction makes $S$ normal and gives an effective
rational dlt different on $S$; its lc centers are precisely the
smaller ambient strata in $S$
\cite{Kollar2013}. The divisor $E_Q$ does not contain $S$. Since
$Q$ is $\Q$-factorial, a Cartier multiple of $E_Q$ restricts to an
effective Cartier divisor on $S$. Adding that restriction to the
different gives an effective rational divisor $B_S$ such that
\begin{equation}
\label{common:subadjunction-stratum}
 K_S+B_S\sim_{\Q}(p|_S)^*\bigl((K_Z+\Theta)|_V\bigr).
\end{equation}
Over the generic point of $V$ this pair is klt: the surplus is
absent there, and a non-klt center of the different dominating
$V$ would be a smaller stratum, contradicting the choice of $S$.

Factor the map through the normalization $V^\nu$ and its Stein
factorization:
\[
 S\xrightarrow{q}V'\xrightarrow{\eta}V^\nu\xrightarrow{\nu}V.
\]
Here $q$ has connected fibers and $\eta$ is finite. Put
$D_V=\nu^*((K_Z+\Theta)|_V)$ and
$D'=\eta^*D_V$. We apply the canonical bundle formula on the
connected-fiber base $V'$, then push the resulting intersection
inequality through the finite map $\eta$. Its ramification will
contribute an effective term. The final comparison is between the
canonical cycle on $V^\nu$ and the canonical divisor on $V^*$.

We verify the klt-trivial-fibration hypotheses for
$q:(S,B_S)\to V'$. The generic pair is sub-klt, and
Equation~\eqref{common:subadjunction-stratum} is the required rational-linear
pullback expression. For the rank condition, take a log resolution
over the generic point of $V'$. Since the boundary on $S$ is
effective and the generic pair is klt, the round-up of its
discrepancy divisor has coefficient zero along every strict
boundary component and has nonnegative coefficients only on
exceptional divisors. Over this dense klt open, its pushforward to the normal $S$ is
$\cO_S$: a rational function allowed poles only over exceptional
centers extends across codimension at least two. Connectedness
in the Stein factorization therefore gives generic rank one for
the required discrepancy round-up pushforward. This verifies the
rank condition, rather than deducing it from connected fibers
alone.

The canonical bundle formula, with Ambro's moduli positivity in
the form recalled by Fujino--Gongyo, now gives
\begin{equation}
\label{common:subadjunction-cbf}
 D'\sim_{\Q}K_{V'}+B_{V'}+M_{V'},
\end{equation}
where $M$ is a b-nef moduli divisor
\cite[Definition~3.1 and Theorem~3.5]{FujinoGongyoModuli2014}.
The definition requires sub-klt over the generic point; it does
not require the pair to be globally sub-lc. Thus the possible
non-lc fibers caused by $E_Q|_S$ are permitted.

The discriminant $B_{V'}$ on this initial base is effective. To
check this at a prime divisor $F\subset V'$, work over its generic
point, where the base is regular. A component of its inverse
image on $S$ has multiplicity $m_F\geq1$ and nonnegative boundary
coefficient $b_F$. Its discrepancy bounds the generic lc threshold
$t_F$ from above by
\[
 t_F\leq\frac{1-b_F}{m_F}\leq1.
\]
Thus the discriminant coefficient $1-t_F$ is nonnegative. If the
fiber is already non-lc, $t_F$ can be negative; this only increases
that coefficient.

Let $P'=\eta^*\nu^*P$. On a higher projective base
$\mu:\widehat V'\to V'$ where the moduli divisor is nef, its
trace satisfies $M_{V'}=\mu_*M_{\widehat V'}$. The projection
formula gives
\[
 M_{V'}\cdot(P')^{f-1}
 =M_{\widehat V'}\cdot(\mu^*P')^{f-1}\geq0.
\]
Only this intersection is asserted to be nonnegative; the trace
$M_{V'}$ need not itself be nef. Likewise
$B_{V'}\cdot(P')^{f-1}\geq0$ by effectivity and nefness.
All intersections with Weil divisors here mean their cycle
intersections with rational Cartier powers, so they are defined
even when the separate canonical or discriminant divisor is not
$\Q$-Cartier.

Choose compatible canonical Weil divisors. Since $\eta$ is finite
and separable, the codimension-one ramification formula, followed
by pushforward, is
\[
 \eta_*K_{V'}=(\deg\eta)K_{V^\nu}+\eta_*R_\eta,
 \qquad R_\eta\geq0.
\]
This identity of Weil cycles uses the discrete valuation rings at
generic prime divisors and does not require $K_{V^\nu}$ to be
$\Q$-Cartier. Intersect Equation~\eqref{common:subadjunction-cbf} with
$(P')^{f-1}$ and discard the nonnegative discriminant and moduli
intersections. Pushforward through $\eta$ and the nonnegative
ramification intersection then give
\[
 (\deg\eta)D_V\cdot(\nu^*P)^{f-1}
 \geq(\deg\eta)K_{V^\nu}\cdot(\nu^*P)^{f-1}.
\]
Finally the pushforward of $K_{V^*}$ to $V^\nu$ is its compatible
canonical cycle. All additional divisors on the resolution have
image of codimension at least two, so their intersection with
$f-1$ pulled-back Cartier powers is zero. Dividing by
$\deg\eta$ proves Equation~\eqref{common:subadjunction-inequality}. When $f=1$
there are no birational exceptional divisors on the normal curve,
and the same calculation applies directly.
\end{proof}

\subsection{Degree and canonical bounds for a base component}

We next apply numerical subadjunction to a boundary built from an
actual linear system. Its local multiplicity controls both the
degree of the center and the coefficient of that boundary.

\begin{lemma}[A base component of high multiplicity]
\label{common:base-component}
Let $Z$ be a projective $\Q$-factorial klt variety of dimension $d$,
let $P$ be a nef rational Cartier divisor, and let $k>0$ be an
integer for which $kP$ is Cartier. Let
$0\ne\mathcal H\subset H^0(Z,\cO_Z(kP))$ be a linear subspace
with proper base locus and base ideal $\mathfrak b$.
Suppose that $V$ is an integral base-locus component of dimension
$f<d$, isolated at its generic point $\eta_V$, and that
$\eta_V\in Z_{\rm sm}$. Write $a=d-f$ and
$\mathfrak m=\mathfrak m_{\cO_{Z,\eta_V}}$. If
\[
 \mathfrak b\cO_{Z,\eta_V}\subset\mathfrak m^h
 \qquad(h\in\Z_{>0}),
\]
then
\begin{equation}
\label{common:base-component-degree}
 P^f\cdot V\leq (k/h)^aP^d.
\end{equation}
There are a rational number $c$ with $0<c\leq ak/h$ and an
effective rational divisor $\Theta\sim_{\Q}cP$ such that
$(Z,\Theta)$ is lc at $\eta_V$ and has an lc place with center
$V$. If $f>0$, on a smooth resolution $V^*$ one consequently has
\begin{equation}
\label{common:base-component-canonical}
 K_{V^*}\cdot P^{f-1}
 \leq (K_Z+cP)|_V\cdot P^{f-1}.
\end{equation}
The degree estimate includes $f=0$; the canonical estimate is
asserted only for positive-dimensional $V$.
\end{lemma}

\begin{proof}
The local ring $R=\cO_{Z,\eta_V}$ is regular of dimension $a$.
Isolation means that $\mathfrak bR$ is $\mathfrak m$-primary.
Choose $a$ general sections of $\mathcal H$. Their local equations
form a system of parameters in $R$, since an $\mathfrak m$-primary
ideal is contained in no smaller prime. Let $\mathfrak q$ be the
parameter ideal they generate. It is contained in $\mathfrak m^h$.
Regularity and the multiplicity comparison give
\[
 \operatorname{length}(R/\mathfrak q)
 =e(\mathfrak q,R)\geq e(\mathfrak m^h,R)=h^a.
\]
This is the generic intersection multiplicity along $V$.

The global intersection can be performed without assuming that
the series has no other base components. Before each cut, discard
the components of the current effective cycle that lie in the
base locus. Before the final cut their dimensions are greater
than $f$, so none can contain $V$: otherwise $V$ would not be a
base-locus component. Discarding them therefore leaves the
calculation at $\eta_V$ unchanged. A general next section meets
the remaining components properly. Nefness of $P$ shows that
discarding effective components can only decrease their total
$P$-degree. After $a$ cuts the surviving effective cycle has
$P$-degree at most $k^aP^d$, and contains $V$ with multiplicity at
least $h^a$. This proves Equation~\eqref{common:base-component-degree}.

Let $\lambda=\operatorname{lct}_{\eta_V}(\mathfrak b)$ be the generic ideal
threshold. It is positive and rational. The valuation from
blowing up the smooth generic center has log discrepancy $a$
and ideal order at least $h$, so
\[
 0<\lambda\leq a/h.
\]
Because $\mathfrak bR$ is primary, a divisor computing this
threshold has center $V$. We realize the ideal threshold by a
divisor rather than apply subadjunction to an ideal formally.
Take a log resolution of $Z$ and $\mathfrak b$, so that the
pulled-back system has fixed divisor $F$ and a basepoint-free
moving part. Choose sufficiently many general members
$D_1,\ldots,D_N$ of $\mathcal H$ and put
\[
 \Theta=\frac{\lambda}{N}(D_1+\cdots+D_N)
 \sim_{\Q}cP,\qquad c=k\lambda.
\]
On the resolution the fixed contribution is $\lambda F$.
The moving contributions have coefficients $\lambda/N<1$ and,
by general choice, meet the fixed resolution divisor with simple
normal crossings on the open in question. Thus $(Z,\Theta)$ is
lc at the generic point of $V$ and has the same lc place there
as the ideal threshold. Global log canonicity is unnecessary.
For $f>0$, Lemma~\ref{common:numerical-subadjunction} gives
Equation~\eqref{common:base-component-canonical}, and
$c=k\lambda\leq ak/h$.
\end{proof}

\subsection{Differentiating the full moving jet kernel}

Fixing a high order at a varying point produces a family of
linear subspaces of one fixed section space. A component common
to two successive base loci acquires high multiplicity, because
parameter derivatives of the higher kernel still belong to the
lower kernel. We first state the parameter-family conclusion. Local restriction
to a fiber is treated separately below. The differentiation argument
is related to the method of Ein--K\"uchle--Lazarsfeld
\cite[Proposition~2.3]{EinKuchleLazarsfeld1995}; the ordinary-power
and primary-ideal details needed here are included.

\begin{lemma}[Moving multiplicity]
\label{common:moving-kernel}
Let $Z$ be an integral projective variety of dimension $d$, and
let $P$ be a nef, big, semiample rational Cartier divisor. Fix an integer $k$
with $kP$ Cartier and positive real numbers
\[
 \tau_0<\tau_1<\cdots<\tau_d,
 \qquad\tau_{i+1}-\tau_i=\delta>0.
\]
For a smooth point $x$, let $\mathcal H_i(x)$ be the entire
subspace of sections of $kP$ vanishing at $x$ to order at least
$\lceil k\tau_i\rceil$. Suppose that for very general $x$ all these
subspaces are nonzero and the base locus of $\mathcal H_0(x)$
has a positive-dimensional component through $x$.
If $k\delta\geq4$, then, after an algebraic parameter extension
and restriction to nonempty opens, there are an irreducible
parameter space $T$, a fixed adjacent pair of levels, and a
family of integral subvarieties $V_t$ common to the two base
loci such that the incidence
\[
 \Gamma=\{(t,z):z\in V_t\}\subset T\times Z
\]
dominates $Z$. Each general $V_t$ has dimension in $[1,d-1]$,
is an isolated component of the higher base locus at its generic
point, and its entire higher-series base ideal is contained
there in
\begin{equation}
\label{common:moving-multiplicity-order}
 \mathcal I_{V_t}^{h},\qquad h=\lceil k\delta/2\rceil,
 \qquad k/h\leq2/\delta.
\end{equation}
These are ordinary local ideal powers in the smooth ambient open.

The incidence base ideal lies in $\mathcal I_\Gamma^h$ on a dense
open of $\Gamma$. If, in addition, $Z$ is $\Q$-factorial and klt,
then for a general member, with $f=\dim V_t$, one has
\begin{equation}\label{common:tracking-bounds}
 \begin{split}
 P^f\cdot V_t&\le(2/\delta)^{d-f}P^d,\\
 K_{V_t^*}P^{f-1}&\le(K_Z+cP)|_{V_t}P^{f-1},
 \qquad 0<c\le2(d-f)/\delta .
 \end{split}
\end{equation}
In particular the weaker bound $c\le2d/\delta$ also holds.
Positivity of $P^f\cdot V_t$ follows if $V_t$ meets the open
where the big semiample contraction is an isomorphism. No such
positivity is asserted for every exceptional subvariety.
\end{lemma}

\begin{proof}
On the smooth marked-point open, evaluation into the finite jet
bundle is a morphism from the constant bundle
$H^0(Z,kP)\otimes\cO$ to a locally free sheaf. Remove the rank-jump
loci for the finitely many levels. Its kernels are then vector
subbundles whose fibers are exactly the full spaces
$\mathcal H_i(x)$. Their base loci are nested in increasing order
as $i$ increases. Starting with the stipulated component through
the mark, follow a component containing it at each later level.
Every such component has dimension between $1$ and $d-1$: the
kernel is nonzero, so its base locus is proper. Among the $d+1$
nested components two successive ones have the same dimension
and hence are equal.

This tracking can be made algebraic for the fixed $k$. Follow the
finitely many components of the geometric generic base loci,
their inclusions, and their incidences with the mark. They are
defined after a finite extension of the generic parameter field.
Spread over a corresponding irreducible parameter variety and
shrink for the required flatness, geometric integrality, and
constant-rank conditions. There are only finitely many choices
of adjacent step; one of them gives a family dominant over the
marked-point open. It gives the asserted $T$ and $\Gamma$.
The evaluation $\Gamma\to Z$ is dominant because the moving mark
belongs to each $V_t$ and its map to $Z$ is dominant.

We now prove ordinary high-order vanishing. Work on a smooth
parameter open, and denote the higher-kernel bundle there by
$\mathcal K$. Evaluation on $T\times Z$ gives a map from the
pullback of $\mathcal K$ to the pullback of $\cO_Z(kP)$.
Let $\mathfrak b_T$ be its base ideal. After trivializing the
target line bundle, a local frame of $\mathcal K$ generates
$\mathfrak b_T$. We show that every frame section lies in
$\mathcal I_\Gamma^h$.
Each frame element has the form
\[
 s(t,z)=\sum_\alpha c_\alpha(t)s_\alpha(z),
\]
where the $s_\alpha$ are a fixed basis of the global space
$H^0(Z,kP)$. Parameter differentiation with $z$ fixed only
differentiates the coefficients. Every derivative, at each fixed
parameter, is consequently a global section of the same line
bundle $kP$.

Write $x(t)$ for the moving mark. In local coordinates $z-x(t)$,
the section starts in order $\lceil k\tau_{i+1}\rceil$. Each parameter
derivative lowers this order by at most one. A derivative of
order $r$ therefore belongs to the lower kernel if
\[
 r\leq\lceil k\tau_{i+1}\rceil-\lceil k\tau_i\rceil.
\]
The right side is at least $k\delta-1$. Since $k\delta\geq4$,
every derivative of order less than $h=\lceil k\delta/2\rceil$ is
allowed. It thus vanishes along $\Gamma$, which lies in the
lower base locus. A change of higher-kernel frame contributes
only parameter coefficients and derivatives of orders no larger
than $r$, so the conclusion holds for any local frame.

Here a first-order normal-space assertion must be upgraded to
all orders. At a general smooth point of $\Gamma$, its projection
to $Z$ is submersive. Thus
\[
 T_{(t,z)}(T\times Z)
 =T_{(t,z)}\Gamma+(T_tT\times\{0\}).
\]
In particular, parameter directions span the normal space to
$\Gamma$. Choose parameter coordinates
$t=(t',t'')$, where $t'$ has length
$a=\codim_{T\times Z}\Gamma$, such that the normal equations have
an invertible Jacobian minor in the $t'$ directions. The analytic
implicit-function theorem writes
\[
 \Gamma:\quad t'=F(t'',z).
\]
Set $u=t'-F(t'',z)$. With $t''$ and $z$ held fixed, differentiation
in $t'$ is exactly differentiation in $u$, to every order. Use a
frame of $kP$ depending only on $z$. The vanishing on $\Gamma$
of all parameter derivatives of order below $h$ says that every
normal Taylor coefficient of degree below $h$ vanishes. Therefore
$s\in\mathcal I_\Gamma^h$ in the ordinary analytic local ring.
Equivalently, in arbitrary relative normal coordinates a product
of $r$ normal vector fields expands into parameter derivatives
of orders at most $r$, so no unaccounted higher-coordinate terms
arise.

The sections and ideals are algebraic. Faithful flatness of the
analytic local ring over the algebraic local ring gives the same
ideal membership algebraically. Applying this to a finite frame
and clearing the finitely many denominators gives containment
on a dense algebraic open of the incidence. Generic smoothness
of $\Gamma\to T$ identifies its scheme fiber there with the
reduced general component $V_t$. Restriction therefore gives
$\mathcal I_\Gamma^h\cO_{\{t\}\times Z}=\mathcal I_{V_t}^h$ locally.
The specialized frame is a basis of the entire higher kernel,
so the containment concerns its whole base ideal.

If $Z$ is $\Q$-factorial and klt, apply
Lemma~\ref{common:base-component} to the higher kernel.
Its generic base ideal is primary because $V_t$ is a base-locus
component, and its generic point is smooth by incidence dominance.
The bounds follow from $k/h\le2/\delta$. On the isomorphism open
of the big semiample contraction the image of $V_t$ has dimension
$f$, so the pullback of an ample class has positive top intersection
on $V_t$. This proves the stated positivity qualification.
\end{proof}

\subsection{Restriction to a suitable fiber}

The restriction step is local and must be separated from the
global positivity and singularity assumptions of subadjunction.

\begin{lemma}[Restriction of an isolated base ideal]
\label{common:fiber-restriction}
In the setting of Lemma~\ref{common:moving-kernel}, let $g:Z\to B$
be a morphism. There is a dense open of the incidence on which
the following holds. Choose $(t,z)$ in that open, let $b=g(z)$,
and suppose that $V_t$ is smooth over the smooth open of its
actual image at $z$. Let $F$ be the reduced component through
$z$ of its scheme fiber. Locally on $Z_b$ at $z$, restriction of
the higher kernel's base ideal has support exactly $F$ and is
contained in $\mathcal I_F^h$.
If $F$ is proper in an integral ambient fiber component containing
it, the restricted series is nonzero there and $F$ is an isolated
base component at its generic point. Application of
Lemma~\ref{common:base-component} on that ambient component
requires all its hypotheses, including projectivity, $\Q$-factoriality,
klt singularities, and smoothness at the generic point of $F$,
separately; they are not conclusions of this restriction lemma.
\end{lemma}

\begin{proof}
Choose the incidence point after removing the other base components,
the locus where the ideal containment is unavailable, and the
failures of generic smoothness for the incidence over its parameter
space and for the component over its actual image. These are proper
closed subsets after shrinking. A prescribed very general condition
on the evaluation can also be imposed because the incidence dominates
$Z$. The component is then chosen through this point, rather than
by selecting a possibly special fiber in advance.

Locally the original base support is exactly $V_t$. Smoothness
over the actual image makes its scheme fiber reduced at the chosen
point. Thus in the ambient fiber
\[
 \mathcal I_{V_t}\mathcal O_{Z_b}=\mathcal I_F,\qquad
 \mathcal I_{V_t}^{h}\mathcal O_{Z_b}=\mathcal I_F^{h}
\]
near that point. Restricting the actual sections extends their
base ideal, which proves both asserted local properties.
No other base component contains the chosen point, so none can
contain the component $F$ through it. If this component is proper
in the chosen integral ambient component, the base support is
proper there, and the restricted series cannot vanish identically.
At its generic point the resulting ideal is primary. This is a
statement about this reduced general fiber and does not assert
transversality or nonzero restriction for arbitrary special fibers.
\end{proof}

\subsection{A fixed branch complement in a finite-type family}

For later applications, the local moving-center estimates are
followed by a global finiteness step. A degree bound alone does
not make a list of branched covers finite. The following statement
records exactly the additional family information that is needed.
In the correspondence case of Proposition~\ref{prop:large-seshadri},
it will turn bounded degree and branch divisors that do not sweep
the base into finitely many normal covers.

\begin{lemma}[Finite covers after excluding sweeping branch divisors]
\label{common:finite-covers}
Let $Y$ be a normal projective complex variety. Let
$\mathcal V\to T$ be a smooth projective family with integral
fibers, over an integral parameter variety, and let
$g:\mathcal V\to Y$ be a morphism whose general fiber maps
$g_t:V_t\to Y$ are dominant and generically finite of degree
at most a fixed integer $e$.
After a finite parameter extension and shrinking, spread the
geometric generic components of the zero divisor of the relative
Jacobian over $Y_{\rm sm}$ to divisors
$\mathcal R_1,\ldots,\mathcal R_N$ in $\mathcal V$.
Retain the notation $\mathcal V,T,g$ after this base change.
Let $I$ be the set of indices $j$ for which the image closure
$\overline{g_t((\mathcal R_j)_t)}$ is a divisor in $Y$ for general $t$.
Assume
\[
 \overline{g(\mathcal R_j)}\ne Y\qquad(j\in I).
\]
Then, after shrinking again, one smooth nonempty open
$Y^\circ\subset Y$ makes every general finite normal Stein cover
of $g_t$ finite \'etale over $Y^\circ$. There are only finitely
many such covers of $Y$, up to isomorphism over $Y$.
The same conclusion holds simultaneously for finitely many
projections. The open and finite lists may depend on all fixed
parameters of the family.
\end{lemma}

\begin{proof}
The relative determinant is nonzero on a dense open because
the general maps are generically finite in characteristic zero.
Its zero divisor has finitely many irreducible components.
Shrink away vertical components and make a finite parameter extension
to define the geometric generic components individually. Their
closures give the divisors $\mathcal R_j$ in the statement.
Shrink further to keep their fiber-image dimensions constant.
These components account for every branch prime of the general
finite Stein cover: a ramification prime there has a strict
transform on $V_t$, because a proper birational morphism to a
normal variety is an isomorphism at each generic divisorial
point. A divisor exceptional over the Stein cover has image of
codimension at least two on that cover and hence on $Y$, so it
cannot hide a branch prime. Components confined to special
parameters disappear upon shrinking.

Put
\[
 B=Y_{\rm sing}\cup\bigcup_{j\in I}\overline{g(\mathcal R_j)}.
\]
Each closed set in this finite union is proper by hypothesis,
so $B\ne Y$. The finite normal covers are unramified in
codimension one over the smooth open $Y^\circ=Y\setminus B$.
Purity of the branch locus makes them finite \'etale there
\cite[Expos\'e~X, Theorem~3.1]{SGA1}.
The topological fundamental group of $Y^\circ$ is finitely
generated. Indeed, triangulate the compact semialgebraic space
$Y$ compatibly with its closed subset $B$
\cite[Theorem~1.10 in the author's notes]{Coste2007}.
After barycentric subdivision the subcomplex for $B$ is full.
On its complement, normalize the sum of the barycentric
coordinates at vertices outside that subcomplex; decreasing
the other coordinates to zero gives a deformation retraction
onto a finite subcomplex. Its fundamental group is finitely
generated. A finitely generated group has only
finitely many permutation representations in each symmetric
group of degree at most $e$. Riemann existence
\cite[Expos\'e~XII, Theorem~5.1]{SGA1} therefore gives
finitely many finite \'etale covers of these degrees. A finite
normal cover of $Y$ is the normalization in the function field
of its restriction to $Y^\circ$, and is determined by that
restriction. Taking the union of the proper image closures for
finitely many projections proves the simultaneous assertion.
\end{proof}
